\section{Related Work}
\label{sec:related}

\paragraph{Prompt injection in LLM agents.}
Prompt injection---embedding malicious instructions in content an LLM processes as
data---was formalized by \citet{perez2022promptinjection} and has since been studied
extensively for tool-using agents~\citep{greshake2023notwhoyoureexpecting,
liu2024promptinjectionattack}. \citet{yi2023benchmarkingpromptinjection} provide a
systematic taxonomy of injection vectors. The dominant mitigation strategy is
instruction-hierarchy enforcement, in which the model is trained to give lower trust
to content from external sources than to system-level instructions~\citep{
openai2024instructhierarchy, anthropic2024instructionhierarchy}. Our work complements this line by examining whether the \emph{clarification interaction
state} amplifies injection susceptibility in coding agents---and finding that for 8B-class
models, the tool-output channel is in fact the dominant risk, motivating channel-specific
defenses.

\paragraph{Clarification in coding agents and HCI.}
The value of clarification-seeking in task-oriented agents has been studied in
dialogue systems~\citep{radlinski2017clarifyingquestions} and, more recently, in
coding agents specifically. \citet{zhang2024clarifyingquestions} show that ambiguity
resolution through targeted questions significantly improves code quality on
underspecified benchmarks. HiL-Bench~\citep{hilbench2024} introduces a formalized
framework for ``blocked information'' scenarios that require human-in-the-loop
interaction to resolve. ClarEval~\citep{clareval2024} proposes metrics for evaluating
clarification quality beyond task completion. Our work introduces the first \emph{security} dimension to this literature: the
clarification channel is a distinct attack surface whose risk profile---relative to
tool-output injection---is model-scale dependent and requires explicit measurement.

\paragraph{Security of coding agents.}
\citet{pearce2022asleepkeyboard} study Copilot-generated code for security vulnerabilities
and find a non-trivial rate of insecure suggestions. \citet{khoury2023securitychatgpt}
extend this to ChatGPT. These works study the agent as a \emph{producer} of vulnerable
code; we study the agent as a \emph{target} for adversarial manipulation through its
interaction protocol. \citet{yang2024swebench} and related agentic coding benchmarks
focus on task-completion correctness; none explicitly models the adversarial clarification
channel.

\paragraph{ASPI report.}
The most directly related work is the concurrent report by~\citet{aspi2026}, which
empirically demonstrates that clarification-state amplifies injection ASR by 15--20$\times$
across two frontier models in general (non-coding) agentic tasks. We build on their
matched design, extend it to the coding domain with concrete attack categories and
machine-checkable predicates, and---crucially---propose and evaluate defenses that
ASPI does not provide.

\paragraph{Agentic safety and oversight.}
Broader work on LLM agent safety includes constitutional AI~\citep{bai2022constitutionalai},
RLHF for agent alignment~\citep{ouyang2022trainingllm}, and process-level reward
modeling~\citep{lightman2023processreward}. The plan-diff gating defense we propose
(Section~\ref{sec:defenses}) is inspired by recent work on authorizing agent
actions~\citep{zhou2024agentauthorization} and draws on principles from
\citet{wu2024agentsafety}.
